% higgs_composite_qg_SM.tex -- Supplemental Material of the Regular Article "Higgs as a Composite, and Quantum Gravity",
% revised hypotheses.  PROJECT spacetimes, folder Revision_SM_Higgs/paper_prd.
% This file is generated from src/higgs_composite_qg_SM.template.tex by scripts/build_paper.py, which inserts every number
% that a program computes, every table of results, every table of checks and every checksum from the files of the folder
% (the bounds of the agreement of two computations rounded up).  Typed in the template are only the parameters of the
% model and of the numerical method, mathematical constants, the quantum numbers that the hypotheses assign, and measured
% values of the Standard Model with their references.  Build:  python3 paper_prd/scripts/build_paper.py
\documentclass[aps,prd,onecolumn,superscriptaddress,nofootinbib,floatfix,longbibliography]{revtex4-2}
\usepackage{amsmath,amssymb}
\usepackage{graphicx}
\usepackage{bm}
\usepackage{longtable}
\usepackage[colorlinks=true,linkcolor=blue,citecolor=blue,urlcolor=blue]{hyperref}
\hypersetup{pdftitle={Supplemental Material for Higgs as a Composite, and Quantum Gravity},pdfauthor={Patrick Lee Nash}}

\newcommand{\LL}{{\rm LL}}
\newcommand{\SF}{{\rm SF}}
\newcommand{\SH}{{\rm SH}}
\newcommand{\tr}{\mathrm{tr}}
\renewcommand{\theequation}{S\arabic{equation}}
\renewcommand{\thetable}{S\arabic{table}}
\renewcommand{\thefigure}{S\arabic{figure}}
\renewcommand{\thesection}{S\arabic{section}}
% the captions of the long tables end their label with a period, like those of the floats
\makeatletter
\def\LT@makecaption#1#2#3{%
  \LT@mcol\LT@cols c{\hbox to\z@{\hss\parbox[t]\LTcapwidth{%
    \sbox\@tempboxa{#1{#2. }#3}%
    \ifdim\wd\@tempboxa>\hsize
      #1{#2. }#3%
    \else
      \hbox to\hsize{\hfil\box\@tempboxa\hfil}%
    \fi
    \endgraf\vskip\baselineskip}%
  \hss}}}
\makeatother

\begin{document}

\title{Supplemental Material for ``Higgs as a Composite, and Quantum Gravity''}

\author{Patrick Lee Nash, Ph.D.}
\affiliation{PROJECT spacetimes}

\date{November 22, 2026}

\begin{abstract}
This Supplemental Material gives the derivations, the numerical methods, the complete tables of results, the checks and
the data behind every statement of the article. It is written to be read on its own: every symbol is defined here
(Sec.~\ref{sec:notation}), and every equation that the article uses is derived or stated here. Equations, tables and
figures of this document carry the prefix S. All numbers are outputs of programs in Rust, checked by independent
programs in Python (numpy and LAPACK) and, for the group theory and the closed forms, by scripts in the Wolfram Language
in exact arithmetic; Sec.~\ref{sec:checks} lists the checks with their results, and Sec.~\ref{sec:data} the files with
their SHA-256 checksums.
\end{abstract}

\maketitle

\tableofcontents

%=====================================================================================================================
\section{Notation and conventions}
\label{sec:notation}

Units: $\hbar=c=1$; the momentum cutoff $\Lambda_\LL=1$ is the unit of energy and of momentum, and $1/\Lambda_\LL$ the unit
of length and of time. The metric is $\eta^\LL={\rm diag}(1,-1,-1,-1)$. The suffix LL marks the objects of the Lorentz
super-loop (the bi-octonion model of Ref.~\cite{Nash2026}), SF those of the superfluid of the earlier hypotheses that the
revised ones keep, and SH those of the revised hypotheses; a superscript $X$ names the method (Sec.~\ref{sec:methods}).
The \emph{second version} of the hypotheses has the two singlets $\chi_\SH$ and $\chi'_\SH$; the \emph{first version},
computed as the comparison, has the pairing term of $\chi_\SH$ only, and $\chi'_\SH$ is a free field. Table~\ref{tab:symbols}
lists the symbols.

\begingroup\small
\begin{longtable}{p{0.30\textwidth}p{0.65\textwidth}}
\caption{\label{tab:symbols}The symbols of the article and of this Supplemental Material.}\\
\hline\hline
symbol & meaning\\ \hline
\endfirsthead
\hline symbol & meaning\\ \hline
\endhead
\hline\hline
\endlastfoot
$\Psi_\LL$; $\chi_\SH$, $\chi'_\SH$ & the isodoublet of Dirac fields (eight components, component $2d+s$: Dirac index $d=0,\dots,3$ of the chiral basis, isospin index $s=0$ up, $s=1$ down); the two singlet Dirac fields (four components each) \\
$m_\LL$, $m^\chi_\SH$, $m^{\chi'}_\SH$ & the masses of the three fields, all $0.2\,\Lambda_\LL$ \\
$\Lambda_\LL$ & the momentum cutoff of the uniform states, the unit of the computation \\
$\gamma^\mu_\LL$, $\beta_\LL=\gamma^0_\LL$, $\alpha^j_\LL$, $\gamma^{\rm ch}_\LL$, $C_\LL$ & the Dirac matrices of the chiral basis, the chirality matrix ($\gamma^5$ of the article) and the charge-conjugation matrix \\
$SU(2)_L$, $U(1)_Y$; $T^{a}$, $\tau^a$ & the electroweak gauge group; its generators ($a=1,2,3$; $\tfrac12\tau^a$ on the isospin of $\Psi_\LL$, zero on the singlets) and the Pauli matrices \\
$Y$, $Q=T^3+Y$ & the hypercharge and the electric charge (Table~\ref{tab:fields}) \\
$g$, $g'$, $\theta_W$ & the couplings of $SU(2)_L$ and $U(1)_Y$ and the weak mixing angle, $\cos\theta_W=M_W/M_Z$ \\
$W^a_\mu$, $B_\mu$; $W^\pm_\mu$, $Z_\mu$ & the gauge fields of $SU(2)_L$ and $U(1)_Y$; the massive weak gauge fields \\
$F^{\rm tot}_\SH$, $R_\SH$ & the particle number (1 on every particle of every field) and the relative number (particles of $\Psi_\LL$ minus particles of the singlets) \\
$SU(2)_R$, $T^a_R$; $SU(2)_V$ & the custodial group, which rotates $(\chi_\SH,\chi'_\SH)$ as a doublet, and its generators; the diagonal subgroup of $SU(2)_L\times SU(2)_R$ \\
$\Pi^s_\SH$, $\Pi'^s_\SH$ & the pair operators (\ref{eq:pairs}) of $\Psi_\LL$ with $\chi_\SH$ and with $\chi'_\SH$ \\
$\mathbf e^{\rm u}$, $\mathbf e^{\rm d}$ & the isospin columns $(1,0)^T$ and $(0,1)^T$ \\
$\Phi^s$, $\Phi'^s$; $\mathcal M_\SH$ & the pair fields, proportional to $\langle\Pi^s_\SH\rangle$ and $\langle\Pi'^s_\SH\rangle$; their $2\times2$ matrix (\ref{eq:M}) \\
$\Pi_\SH[\mathsf V_\SH]$, $\Xi_\SH[\mathsf V_\SH]$ & a pair and a particle-antiparticle bilinear with the coefficient matrix $\mathsf V_\SH$ (\ref{eq:bilinears}) \\
$G^{\rm S}_\SF$, $G^{\rm V}_\SF$, $G_\SH$ & the couplings of the scalar, the vector and the pairing contact terms ($-8$, $0.8$, $100.8$ in units of $\Lambda_\LL^{-2}$) \\
$V^{(0)}$, $V^{(1)}$ & the contact potential without the pairing terms, and with them: the potential of the article \\
$I_1(M)$; $G^{\rm crit}_\SH$ & the integral (\ref{eq:I1}); the critical coupling $4\pi^2/I_1(m_\LL)$ of the pairing \\
$\mathsf h(\vec q)$ & the free Dirac Hamiltonian of the sixteen components at the momentum $\vec q$ \\
$\mathsf N$, $\mathsf P$; $\mathsf H(\vec q)$ & the normal and the pair self-energies of a uniform state; its Bogoliubov-de Gennes (BdG) matrix (\ref{eq:BdG}) \\
$\varrho$, $\kappa$ & the local normal matrix, measured from the normal vacuum, and the local pair matrix \\
$\mathcal E^X$ & the energy per volume of a uniform state in method $X$, measured from the normal vacuum \\
$X$; $s_{\rm S}$, $s_{\rm V}$ & a method: HS (1,1), HS (0,1), HS (0,0), KS or CJT; the shares of the scalar and the vector term decoupled in density channels in HS \\
$M^{\rm u}$, $M^{\rm d}$, $M^\chi$, $M^{\chi'}$ & the Dirac masses of the four species in a state: the bare mass plus a quarter of the trace of $\mathsf N$ with the matrix $\beta_\LL$ of the species \\
$\varsigma^{(1)}$, $\varsigma^{(2)}$ & the singular values of the $2\times2$ matrix of the coefficients of $(\mathsf K^s)^\dagger$ and $(\mathsf K'^s)^\dagger$ in $\mathsf P$, the larger first \\
$\mathsf K^s$, $\mathsf K'^s$ & the $16\times16$ pair matrices of $\Pi^s_\SH$ and $\Pi'^s_\SH$ \\
$\mathsf P^{\rm vac}$; $d$, $d^0_\SH$ & the direction (\ref{eq:Pvac}) of the vacuum in the pair self-energy; the coefficient of a pair self-energy along it, and its value in the ground state \\
$\varphi_\SH$, $\varphi^0_\SH$ & the pair amplitude $|\langle\Pi^{\rm d}_\SH\rangle|$ of a state, and its value in the ground state \\
$E^{X,{\rm gap}}$ & the gap of a ground state: the lowest energy of a quasi-particle \\
$\lambda^{\rm H}$, $\lambda^{\rm S}$, $\lambda^{\rm T}$ & the largest eigenvalues of the static response of the normal vacuum in the doublet-singlet pairs, the scalar channel and the triplet pairs of $\Psi_\LL$ \\
$V^{\rm eff}_\SH$; $V^{\rm eff\,\prime\prime}_\SH$ & the effective potential along the direction of the vacuum; its second derivative \\
$Q_a$, $\mu^a$, $\chi^{ab}$ & the eight generators of the global group as $16\times16$ matrices, the chemical potential of $Q_a$, and the susceptibilities (\ref{eq:chi}); $\chi^{11}$, $\chi^{33}$, $\chi^{YY}$, $\chi^{3Y}$, $\chi^{FF}$ those of $T^1$, $T^3$, $Y$, the mixed one of $T^3$ and $Y$, and that of $F^{\rm tot}_\SH$ \\
$Z_\SH$; $\mathsf M^{\rm g}$ & the kinetic coefficient of the pair fields; the mass matrix (\ref{eq:Mg}) of the gauge fields \\
$v_\SH$, $\rho_\SH$; $M_{W,\SH}$, $M_{Z,\SH}$ & the scale of the breaking and the ratio $M_W^2/(M_Z^2\cos^2\theta_W)$ of the superfluid; its masses of the $W$ and the $Z$ \\
$v$, $m_H$ & the vacuum value of the Higgs field of the Standard Model and the mass of the Higgs boson \\
$M^{\rm EFT}_\SH$; $\omega^{\rm H}_\SH$ & the frequency (\ref{eq:MEFT}) of the amplitude mode from the effective Lagrangian; its frequency from the linear response \\
$M_{\rm eff}$ & the energy $(m_\LL^2+(d^0_\SH)^2/4)^{1/2}$ of a paired quasi-particle at rest in the pairing-only methods \\
$\bm\chi(\omega)$, $\mathsf F^X$ & the linear response of a uniform state at the frequency $\omega$, and the kernel of the contact terms of method $X$ \\
$\lambda^{\rm amp}(\omega)$, $\lambda^{\rm ph}(\omega)$ & the eigenvalues (\ref{eq:lam}) of the response of the pair amplitudes in HS (0,0) along the amplitude and the phase of the vacuum direction \\
$\epsilon^{\rm kick}$ & the relative size of the displacement of the pair self-energy that starts a run in time \\
$\omega^{\rm pk}$, $\Gamma^{\rm pk}$, $\Gamma^{\rm win}$ & the frequency of a peak of the spectrum of a run, its half-width at half maximum, and the half-width that the finite length of the run alone gives \\
$\mathsf O^\dagger\mathsf O-\mathsf 1$ & the departure of the quasi-particle vectors of a run from orthonormality \\
$q$, $\vec q$, $M$, $E$ & a momentum, its vector, a mass and an energy (integration variables) \\
\end{longtable}
\endgroup

%=====================================================================================================================
\section{The model}
\label{sec:model}

\subsection{Fields and charges}

There are three Dirac fields of $SL(2,\mathbb C)$, all singlets of colour and all vector-like. Table~\ref{tab:fields} gives
their charges. The particle number $F^{\rm tot}_\SH$ and the relative number $R_\SH$ are global symmetries; neither is
gauged. On $\Psi_\LL$ and $\chi_\SH$ the hypercharge is a quarter of the fermion number $F_\LL$ of the earlier hypotheses,
$Y=\frac14F_\LL$, so the gauge field of $U(1)^{\rm F}$ of those hypotheses is $B_\mu$ with the coupling $g'/4$; on
$\chi'_\SH$, $4Y=-3F^{\rm tot}_\SH$. The isospin group $SU(2)_{\rm v}$ of the earlier hypotheses is $SU(2)_L$.

\begin{table}[h]
\caption{\label{tab:fields}The fields and their charges.}
\begin{ruledtabular}
\begin{tabular}{lcccccc}
field & $SU(2)_L$ & $Y$ & $Q$ & $F^{\rm tot}_\SH$ & $R_\SH$ & mass\\
\hline
$\Psi_\LL$, up & $\mathbf 2$ & $\frac14$ & $\frac34$ & 1 & $+1$ & $m_\LL$\\
$\Psi_\LL$, down & $\mathbf 2$ & $\frac14$ & $-\frac14$ & 1 & $+1$ & $m_\LL$\\
$\chi_\SH$ & $\mathbf 1$ & $\frac14$ & $\frac14$ & 1 & $-1$ & $m^\chi_\SH$\\
$\chi'_\SH$ & $\mathbf 1$ & $-\frac34$ & $-\frac34$ & 1 & $-1$ & $m^{\chi'}_\SH$\\
\end{tabular}
\end{ruledtabular}
\end{table}

The covariant derivative is
\begin{equation}
D_\mu=\partial_\mu+\Gamma^\LL_\mu+ig\,\tfrac{\tau^a}{2}W^a_\mu+ig'\,YB_\mu ,
\label{eq:cov}
\end{equation}
with the Lorentz connection $\Gamma^\LL_\mu$, a pure gauge on flat space-time, set to zero. The couplings are those of the
Standard Model, from the measured Fermi constant and masses of the weak bosons~\cite{PDG2026}:
$v=(\sqrt2G_F)^{-1/2}={{V_GEV}}$~GeV, $g=2M_W/v={{G}}$, $\cos\theta_W=M_W/M_Z$ and $g'=g\tan\theta_W={{GP}}$.
Because every field is vector-like, every trace of the gauge anomalies vanishes (checked:
Sec.~\ref{sec:checks}).

\subsection{Bilinears and channels}

The bilinears of spin zero of two fields are the pairs and the particle-antiparticle bilinears,
\begin{equation}
\Pi_\SH[\mathsf V_\SH]=\Psi^T_\LL\mathsf V_\SH\chi_\SH,\qquad \Xi_\SH[\mathsf V_\SH]=\Psi^\dagger_\LL\mathsf V_\SH\chi_\SH ,
\label{eq:bilinears}
\end{equation}
and likewise for every other choice of the two fields. A bilinear is a Lorentz scalar when the coefficient matrix obeys
$\mathsf s(\mathbf u)^T\mathsf V+\mathsf V\mathsf s(\mathbf u)=0$ (pairs) or
$\mathsf s(\mathbf u)^\dagger\mathsf V+\mathsf V\mathsf s(\mathbf u)=0$ (particle-antiparticle) for the six generators
$\mathbf u$ of $SL(2,\mathbb C)$, with $\mathsf s(\mathbf u)$ their action on the Dirac index. For two particles of one field
only the antisymmetric part of $\mathsf V$ gives an operator, because the components anticommute. The Lorentz-invariant
bilinears of one kind and one parity form a channel; a generator of the internal group acts on a channel, in a basis
orthonormal in $\tr(\mathsf V^\dagger\mathsf V')$, as a Hermitian matrix, whose eigenvalues are the charges of the channel.
Table~\ref{tab:channels} gives the pair channels of positive parity. The two Lorentz-invariant matrices of a pair of Dirac
spinors, $C_\LL$ and $C_\LL\gamma^{\rm ch}_\LL$, are antisymmetric, so for two members of $\Psi_\LL$ the isospin factor must be
symmetric:
\begin{equation}
\mathbf 2\otimes\mathbf 2=\mathbf 1_{\rm antisymmetric}\oplus\mathbf 3_{\rm symmetric}.
\label{eq:2x2}
\end{equation}
The pairs of two members of $\Psi_\LL$ are therefore a triplet with $Y=\frac12$ and the charges $\frac32$, $\frac12$,
$-\frac12$, without a neutral component; only the pairs of $\Psi_\LL$ with a singlet are doublets with $|Y|=\frac12$ and a
neutral component. Every pair without $\Psi_\LL$ and a singlet has $|Q|\ge{{MIN_ABS_Q}}$ (checked in exact arithmetic).

{{T_CHANNELS}}

\subsection{The interaction and the custodial symmetry}

The pair operators of the attraction are
\begin{equation}
\Pi^s_\SH=\tfrac12\Pi_\SH[C_\LL\gamma^{\rm ch}_\LL\otimes\mathbf e^s],\qquad \Pi'^s_\SH=\tfrac12\Pi_\SH[C_\LL\gamma^{\rm ch}_\LL\otimes\mathbf e^s]\big|_{\chi_\SH\to\chi'_\SH},
\label{eq:pairs}
\end{equation}
$s={\rm u},{\rm d}$, and the Hamiltonian density is
\begin{multline}
\mathcal H=\Psi^\dagger\mathsf h(-i\nabla)\Psi+\tfrac12G^{\rm S}_\SF{:}(\bar\Psi_\LL\Psi_\LL)^2{:}
+\tfrac12G^{\rm V}_\SF{:}\Big[(\Psi^\dagger_\LL\Psi_\LL)^2-\sum_j\big(\Psi^\dagger_\LL(\alpha^j_\LL\otimes\mathsf 1_2)\Psi_\LL\big)^2\Big]{:}\\
-G_\SH\sum_{s={\rm u},{\rm d}}\big[{:}(\Pi^s_\SH)^\dagger\Pi^s_\SH{:}+{:}(\Pi'^s_\SH)^\dagger\Pi'^s_\SH{:}\big],
\label{eq:H}
\end{multline}
where $\Psi=(\Psi_\LL,\chi_\SH,\chi'_\SH)$ has sixteen components, $\mathsf h(\vec q)$ is the free Dirac Hamiltonian,
block-diagonal with $(\alpha^j_\LL q^j+\beta_\LL m_\LL)\otimes\mathsf 1_2$ on $\Psi_\LL$ and $\alpha^j_\LL q^j+\beta_\LL m^\chi_\SH$ and
$\alpha^j_\LL q^j+\beta_\LL m^{\chi'}_\SH$ on the two singlets, the momenta obey $|\vec q|<\Lambda_\LL$, and the colons order the operators so that their expectation value in the
normal vacuum vanishes. The couplings are $G^{\rm S}_\SF=-8$, $G^{\rm V}_\SF=0.8$, $G_\SH=100.8$ (units of $\Lambda_\LL^{-2}$).
In the first version the terms of $\Pi'^s_\SH$ are absent.

Without $B_\mu$, the generators of $SU(2)_L$ (on $\Psi_\LL$), of $SU(2)_R$ ($\frac12\tau^a$ on the pair
$(\chi_\SH,\chi'_\SH)$), $R_\SH$ and $F^{\rm tot}_\SH$ commute with (\ref{eq:H}) when the singlets have equal masses; the
hypercharge is
\begin{equation}
Y=T^3_R+\tfrac14R_\SH ,
\label{eq:Y}
\end{equation}
so $g'$ breaks $SU(2)_R$ to $T^3_R$. Table~\ref{tab:custodial} gives the pair channels of $\Psi_\LL$ with the pair of singlets
under $SU(2)_L\times SU(2)_R$. The pair fields form the matrix
\begin{equation}
\mathcal M_\SH=\begin{pmatrix}\Phi^{\rm u}&\Phi'^{\rm u}\\\Phi^{\rm d}&\Phi'^{\rm d}\end{pmatrix}\to L\,\mathcal M_\SH R^T ,
\label{eq:M}
\end{equation}
a bidoublet: its columns are a doublet with $Y=\frac12$ (charges $+1$, 0) and a doublet with $Y=-\frac12$ (charges 0,
$-1$), whose conjugate $i\tau^2(\Phi')^*$ is a second doublet with $Y=\frac12$. A vacuum with only the neutral components,
of equal lengths, is brought by the symmetries to the form $\mathcal M_\SH\propto i\tau^2$, which is unchanged for $L=R$
because $L\,i\tau^2L^T=i\tau^2\det L$. It keeps $SU(2)_V$, $R_\SH$ and the electric charge
$Q=T^3+T^3_R+\frac14R_\SH$, and breaks three generators of $SU(2)_L\times U(1)_Y$ and $F^{\rm tot}_\SH$.

{{T_CUSTODIAL}}

%=====================================================================================================================
\section{The five mean-field methods}
\label{sec:methods}

\subsection{Quasi-free states}

A uniform quasi-free state is the ground state of a quadratic Hamiltonian with constant self-energies, a Hermitian
$16\times16$ matrix $\mathsf N$ and an antisymmetric one $\mathsf P$. At the momentum $\vec q$ its BdG matrix is
\begin{equation}
\mathsf H(\vec q)=\begin{pmatrix}\mathsf h(\vec q)+\mathsf N&\mathsf P\\\mathsf P^\dagger&-\big(\mathsf h(-\vec q)+\mathsf N\big)^T\end{pmatrix},
\label{eq:BdG}
\end{equation}
and the state fills its negative eigenvalues. Its local matrices are the integrals over the momenta of the normal
correlation matrix, measured from the normal vacuum, and of the pair correlation matrix, $\varrho$ and $\kappa$. The energy
per volume $\mathcal E^X$ is the kinetic part plus the interaction energy that method $X$ assigns to the local matrices.
The pair self-energy of the vacuum direction is
\begin{equation}
\mathsf P=d\,\mathsf P^{\rm vac},\qquad \mathsf P^{\rm vac}=(\mathsf K'^{\rm u})^\dagger-(\mathsf K^{\rm d})^\dagger\ \ \text{(second version)},\qquad
\mathsf P^{\rm vac}=(\mathsf K^{\rm d})^\dagger\ \ \text{(first version)},
\label{eq:Pvac}
\end{equation}
with $d$ real after a rotation by the particle number, and the pair amplitude is
$\varphi_\SH=|\langle\Pi^{\rm d}_\SH\rangle|=\frac12|\tr(\mathsf K^{\rm d}\kappa^T)|$.

\subsection{The methods}

A method keeps the three Wick contractions of every contact term, the direct, the exchange and the pair contraction,
with its own weights:
\begin{itemize}
\item HS $(s_{\rm S},s_{\rm V})$~\cite{Hubbard1959,Stratonovich1957} rewrites each four-fermion term as the exchange of an
auxiliary field at the saddle point; it keeps the direct contraction of the shares $s_{\rm S}$, $s_{\rm V}$ of the scalar and
the vector term decoupled in density channels and the pair contraction of the rest; the pairing terms are decoupled in
their own pair channels; no exchange contraction is kept.
\item KS~\cite{HohenbergKohn1964,KohnSham1965} keeps all three contractions, the exchange evaluated with the projection of
$\varrho$ onto sixteen local densities: the number and the scalar densities of $\Psi_\LL$ as $2\times2$ matrices in the
isospin space, and those of the pair of singlets as $2\times2$ matrices in the space of the two singlets.
\item CJT~\cite{CJT1974} keeps all three contractions with the weight 1, the two-loop effective action of the propagator
truncated at the contact terms.
\end{itemize}
The self-energies are the derivatives of the interaction energy with respect to the local matrices, and a stationary
state is a state whose self-energies are those that its own local matrices produce. In HS (0,1) and HS (0,0) only the
pairing terms act on the states of this paper, and their normal self-energy vanishes: these are the pairing-only methods.

\emph{KS and CJT coincide on the states of this paper.} If $\varrho$ lies in the span of the sixteen Kohn-Sham densities,
the two methods differ only in the projection of the exchange self-energies, and those lie in the span again: for the
scalar and the vector terms because $\beta_\LL^2=(\alpha^j_\LL)^2=\mathsf 1_4$ and $\alpha^j_\LL$ anticommutes with $\beta_\LL$;
for the pairing terms because $\mathsf K^\dagger\varrho^T\mathsf K$ maps a number or scalar density of one singlet to one of
$\Psi_\LL$ times $\mathbf e^s(\mathbf e^s)^T$, and a density of $\Psi_\LL$ to one of a singlet; and no map produces a block
between $\Psi_\LL$ and a singlet. The normal vacuum and every ground state satisfy the hypothesis, because their symmetries
(the rotations, the parity, $R_\SH$, $Q$ and, in the second version, $SU(2)_V$) confine $\varrho$ to the span. KS and CJT
therefore give the same energies, self-energies, susceptibilities and dynamics on them, which the programs confirm to the
rounding. The two methods do differ in the response to changes that leave the span, the mixings of the fields. The
normal blocks between $\Psi_\LL$ and a singlet, and between the two singlets, are reached only by the exchange contraction
of the pairing terms, so HS, which keeps no exchange contraction, has none; in KS the blocks between $\Psi_\LL$ and a
singlet are absent, because the Kohn-Sham densities have no such block, and the block between the singlets is projected
to zero; CJT keeps the exchange contraction, and the largest eigenvalue of its static response is {{MIX_CJT}} in the
blocks between $\Psi_\LL$ and the singlets and {{MIX_SS_CJT}} in the block between the singlets (Table~\ref{tab:normal}).

\subsection{The numerical engine}

The states are computed on 24 Gauss-Legendre nodes of $|\vec q|$ on $[0,\Lambda_\LL]$ times the six directions
$\pm\hat x$, $\pm\hat y$, $\pm\hat z$, by linear mixing of the self-energies (weight 0.6) to a residual below
$10^{-13}$. A calculation on 48 nodes gives the same $v_\SH$ and amplitude mode to {{GRID_DEV}}. The program in Rust
computes every result; the program in Python repeats the key results with its own engine (numpy, vectorized over the
momenta) and by other methods; the scripts in the Wolfram Language derive the group theory and the closed forms in exact
arithmetic or with 30 to 40 digits, and a notebook repeats them.

%=====================================================================================================================
\section{The normal vacuum and the critical coupling}
\label{sec:normal}

In the normal vacuum ($\mathsf N=\mathsf P=0$) the static response to a pair self-energy in the doublet-singlet channel of
the pairing-only methods is $\lambda^{\rm H}=G_\SH I_1(m_\LL)/(4\pi^2)$, with
\begin{equation}
I_1(M)=\int_0^{\Lambda_\LL}\frac{q^2\,dq}{(q^2+M^2)^{1/2}}=\tfrac12\Big[\Lambda_\LL(\Lambda_\LL^2+M^2)^{1/2}-M^2\,{\rm arsinh}\frac{\Lambda_\LL}{M}\Big],
\label{eq:I1}
\end{equation}
so the critical coupling is $G^{\rm crit}_\SH=4\pi^2/I_1(m_\LL)={{GCRIT}}$ and $\lambda^{\rm H}=G_\SH/G^{\rm crit}_\SH={{LAMH}}$.
The scalar and the vector terms do not act on this channel in the normal vacuum, so $\lambda^{\rm H}$ is the same in every
method. Table~\ref{tab:normal} gives the largest eigenvalues of the static response in every channel: only the
doublet-singlet pairs exceed 1, with the multiplicity of the four complex pair fields (eight real directions; four in the
first version). Without the pairing terms (the potential $V^{(0)}$) every start of the iteration ends in the normal vacuum.

{{T_NORMAL}}

%=====================================================================================================================
\section{The stationary states}
\label{sec:states}

The iteration was started from nineteen self-energies: the normal vacuum, pair self-energies along each doublet, along
the custodial and other directions of the bidoublet, along the pairs of one field, shifts of the masses, polarizations of
the isospin and of the singlets, and mixings. Table~\ref{tab:states} gives the stationary states that they reach and the
largest eigenvalue of their static response in the doublet-singlet pairs. With the pairing terms every state other than
the custodial superfluid has an eigenvalue above 1: a small pair amplitude grows on it. The \emph{scalar condensate} of
KS and CJT of the second version, reached from the start that shifts every mass by $-1$, has no pairs and no mixing, and
the Dirac masses of all four species are negative; it occurs only where the exchange contractions couple the scalar
density of $\Psi_\LL$ with those of both singlets.

{{T_STATES}}

Table~\ref{tab:ground} gives the ground states: in the second version the custodial superfluid
($\varsigma^{(1)}=\varsigma^{(2)}$, the matrix of the pair coefficients proportional to $i\tau^2$), in the first the
condensate of $\Phi$ alone ($\varsigma^{(2)}=0$). Table~\ref{tab:sfmasses} gives their Dirac masses. Table~\ref{tab:response}
gives their static response: exactly four eigenvalues equal to 1 in the second version (three gauged generators and the
particle number) and all others below 1. The condensate of one doublet, stationary in the second version, has an
eigenvalue {{ONE_DOUBLET_RANGE}} in the direction of the second doublet.

{{T_GROUND}}

{{T_SFMASSES}}

{{T_RESPONSE}}

%=====================================================================================================================
\section{The effective potential}
\label{sec:potential}

The effective potential $V^{\rm eff}_\SH(\varphi_\SH)$ is the lowest energy per volume of the uniform states with the pair
amplitude $\varphi_\SH$ along the direction of the vacuum. Minimizing the energy at a fixed pair amplitude adds a Lagrange
multiplier times $\Pi^{\rm d}_\SH+(\Pi^{\rm d}_\SH)^\dagger$ (and the same for $\Pi'^{\rm u}_\SH$ in the second version) to the
Hamiltonian, so the constrained state is the ground state of (\ref{eq:BdG}) with $\mathsf P=d\,\mathsf P^{\rm vac}$ for some
$d$ and the normal self-energy that its own local matrices produce. The program fixes $d$, iterates $\mathsf N$ to
self-consistency, and records the energy, the pair amplitude and the Dirac masses, for 33 values
$d=0,0.05\,d^0_\SH,\dots,1.6\,d^0_\SH$, each state started from the normal self-energy of its neighbour. Because the
self-consistent $\mathsf N$ need not be unique at fixed $d$, the continuation is made from both ends: from the ground state
downward to $d=0$ and upward (the branch of the ground state), and from the normal vacuum upward. Where the two reach
different states at the same $d$, the lower branch at the same $\varphi_\SH$ is $V^{\rm eff}_\SH$, and the other is followed
with the step $0.005\,d^0_\SH$ to its end. The derivatives at the minimum follow from central differences in $d$ with the
steps 0.01 and 0.02 combined by Richardson extrapolation, by the parametric formulas
$V'=\mathcal E_d/\varphi_d$ and $V''=(\mathcal E_{dd}-V'\varphi_{dd})/\varphi_d^2$.

In every method and both versions $V^{\rm eff}_\SH$ is a Mexican hat with its minimum at $d^0_\SH$ (Table~\ref{tab:minimum},
Fig.~\ref{fig:potential}). Its top (Table~\ref{tab:top}) is the normal vacuum, except in KS and CJT of the second version,
where the branch of the ground state ends at $d=0$ in the scalar condensate, at
$V^{\rm eff}_\SH(0)={{TOP_RATIO}}\,|V^{\rm eff}_\SH(\varphi^0_\SH)|$; along it the masses of $\Psi_\LL$ stay negative. There the
continuation from the normal vacuum follows a second branch, with positive masses of $\Psi_\LL$, which lies above the
effective potential at the same pair amplitude by {{SB_GAP_RANGE}} times the depth. The continuation with the fine step
follows it up to $d={{SB_LAST}}\,d^0_\SH$, where the number of steps of the iteration has risen steeply (to {{SB_LAST_IT}}),
as near a fold where two branches of stationary states meet and end; at ${{SB_BEYOND}}\,d^0_\SH$ it reaches the branch of
the ground state. The first form of this computation followed only the continuation from the normal vacuum and stated a top
of zero in every method; for KS and CJT of the second version that was incorrect, and the present computation replaces
it. The minimum, its curvature and every later result are computed at the ground state and are unchanged.

{{T_MINIMUM}}

{{T_TOP}}

\begin{figure}[h]
\includegraphics[width=0.9\textwidth]{figures/fig_sh_hg_potential.pdf}
\caption{\label{fig:potential}The effective potential divided by the depth of its minimum, against the pair amplitude,
for the five methods: the second version (left) and the first (right). KS and CJT coincide. The dashed line with circles
is the second branch of KS and CJT of the second version.}
\end{figure}

%=====================================================================================================================
\section{The susceptibilities and the masses of the gauge bosons}
\label{sec:gauge}

\subsection{Chemical potentials as the time components of the gauge fields}

A uniform, static gauge field couples to the charge densities. Its time component $A^a_0$ acts like a chemical potential
$\mu^a=g_aA^a_0$ of the charge $Q_a$. In the effective Lagrangian of the pair fields, $Z_\SH|D_\mu\Phi|^2$ gives for a
time component the term $Z_\SH\mu^2|Q\Phi|^2$, which lowers the energy of the superfluid by the same amount: the
susceptibilities
\begin{equation}
\chi^{ab}=\frac{\partial\langle Q_a\rangle}{\partial\mu^b}=-\frac{\partial^2\mathcal E}{\partial\mu^a\partial\mu^b}
\label{eq:chi}
\end{equation}
of the ground state of $H-\sum_b\mu^bQ_b$ are the coefficients of the mass term. The eight generators are $T^1$, $T^2$,
$T^3$, $T^1_R$, $T^2_R$, $T^3_R$, $R_\SH$ and $F^{\rm tot}_\SH$; the susceptibilities of the hypercharge
$Y=T^3_R+\tfrac14R_\SH$ are the corresponding combinations of theirs, and the mass matrix of $(W^1_\mu,W^2_\mu,W^3_\mu,B_\mu)$ is
\begin{equation}
(\mathsf M^{\rm g})^2_{ab}=g_ag_b\,\chi^{ab},\qquad a,b=1,2,3,Y,\qquad g_1=g_2=g_3=g,\quad g_Y=g',
\label{eq:Mg}
\end{equation}
where the index $a=1,2,3$ stands for $T^a$ and the index $Y$ for the hypercharge.
The space components of a uniform gauge field would also give it, but with the sharp cutoff of the momenta the normal
vacuum answers them with a current, an artifact that the time components avoid: the normal vacuum has a gap and does not
respond to a chemical potential at all. The program computes the ground states of $H-\mu^bQ_b$ for $\mu^b=\pm h$, $\pm2h$
($h=0.005$) and differentiates the charge densities $\langle Q_a\rangle=\tr(Q_a\varrho)$ by central differences with
Richardson extrapolation; by the theorem of Hellmann and Feynman $\partial\mathcal E/\partial\mu^a=-\langle Q_a\rangle$, so
$\chi$ must be symmetric. The Python program computes the same susceptibilities from the second derivatives of the
energy, a second method.

\subsection{Results}

Table~\ref{tab:chi} gives the susceptibilities. In every method they are symmetric, the row of $Q=T^3+Y$ vanishes (the
photon is massless), and the matrix of the eight generators has rank four in the second version; in the first, where of
$SU(2)_R$ only $T^3_R$ is a symmetry, the matrix of $T^1$, $T^2$, $T^3$, $R_\SH$ and $F^{\rm tot}_\SH$ has rank three. In
the second version the custodial structure $\chi(T^a)=\chi(T^a_R)=-\chi(T^a,T^a_R)$ holds, so $\chi^{YY}=\chi^{33}=-\chi^{3Y}$
and the mass matrix (\ref{eq:Mg}) has the eigenvalues $0$, $g^2\chi^{11}$ (twice) and $(g^2+g'^2)\chi^{33}$:
\begin{equation}
M_{W,\SH}=\tfrac12g\,v_\SH,\quad M_{Z,\SH}=\tfrac12(g^2+g'^2)^{1/2}v_\SH,\quad v_\SH=2(\chi^{11})^{1/2},\quad
\rho_\SH=\frac{\chi^{11}}{\chi^{33}}=1 .
\label{eq:rho}
\end{equation}
In the first version $\rho_\SH$ is {{RHO1_RANGE}}, which the measured $M_W^2/(M_Z^2\cos^2\hat\theta_W)={{RHO_MEAS}}$ (with the
$\overline{\rm MS}$ angle $\sin^2\hat\theta_W(M_Z)=0.23122$~\cite{PDG2026}) excludes. Table~\ref{tab:scale} gives the scale
and the masses in GeV, with $v_\SH=v$, and Table~\ref{tab:gapgev} the gap in GeV: the lightest fermionic excitations, of
the charges $\frac34$ and $-\frac14$.

{{T_CHI}}

{{T_SCALE}}

{{T_GAPGEV}}

%=====================================================================================================================
\section{The amplitude mode}
\label{sec:amplitude}

\subsection{The effective Lagrangian}

At long wavelengths the pair fields obey $Z_\SH(|\partial\Phi^{\rm d}|^2+|\partial\Phi'^{\rm u}|^2)-V^{\rm eff}_\SH$ (second
version). A chemical potential $\mu$ of $T^3$, under which $\Phi^{\rm d}$ has the charge $-\frac12$ and $\Phi'^{\rm u}$ the
charge $+\frac12$, lowers the energy by $Z_\SH\mu^2(\frac14+\frac14)(\varphi^0_\SH)^2$, so $\chi^{33}=Z_\SH(\varphi^0_\SH)^2$.
When both lengths oscillate together, $\varphi_\SH=\varphi^0_\SH+\eta$, the kinetic term is $2Z_\SH\dot\eta^2$ and the potential
$\frac12V^{\rm eff\,\prime\prime}_\SH\eta^2$, so
\begin{equation}
(M^{\rm EFT}_\SH)^2=\frac{V^{\rm eff\,\prime\prime}_\SH(\varphi^0_\SH)\,(\varphi^0_\SH)^2}{4\chi^{33}} .
\label{eq:MEFT}
\end{equation}
The same formula holds in the first version, with $\chi^{33}=\frac12Z_\SH(\varphi^0_\SH)^2$ and the kinetic term
$Z_\SH\dot\eta^2$. Formula (\ref{eq:MEFT}) is the leading order of an expansion in the frequency; Table~\ref{tab:amp} shows that
the mode lies close to the energy $2E^{X,{\rm gap}}$ that breaks a pair, where that expansion fails.

{{T_AMP}}

\subsection{The linear response}

A uniform fluctuation with the time factor $e^{-i\omega t}$ changes the self-energies by complex amplitudes $\delta\mathsf N$,
$\delta\mathsf P$ and $\delta\bar{\mathsf P}$ (independent at a real frequency). In the eigenbasis $(E_k,|k\rangle)$ of
(\ref{eq:BdG}) at each momentum the change of the projector onto the negative eigenvalues is
\begin{equation}
\delta R_{kl}=\delta H_{kl}\,\frac{f_l-f_k}{\omega-E_k+E_l},\qquad f_k=1\ \text{for}\ E_k<0,\ 0\ \text{otherwise},
\label{eq:dR}
\end{equation}
summed over the momenta; this is the response $\bm\chi(\omega)$ of the local matrices. The contact terms map the changes of
the local matrices back to the self-energies by the kernel $\mathsf F^X$, and a collective mode is a frequency at which
$1-\bm\chi(\omega)\mathsf F^X$ is singular. The program computes it in the sector of relative number 0 (the normal blocks of
$\Psi_\LL$ and of the pair of singlets and the pairs of $\Psi_\LL$ with the singlets, closed under the symmetries of the
state), in HS (0,1) and HS (0,0), where the mode lies below the threshold $2M_{\rm eff}$ of two paired quasi-particles. It
searches the frequencies below the threshold for zeros of the smallest singular value, finds the null space at each zero,
and projects onto it the amplitude of the vacuum direction; the amplitude mode is the zero whose null space contains it,
and the dimension of the null space is its multiplicity. Table~\ref{tab:rpa} gives the results: at $\omega=0$ exactly
four singular values vanish in the second version (three in the first), the Goldstone modes, and below the threshold the
amplitude mode lies at $\omega^{\rm H}_\SH=d^0_\SH$ to {{OMEGA_D0_DEV}}.

{{T_RPA}}

\subsection{The closed form of the pairing-only methods}

In HS (0,0) the response of the pair amplitudes takes the form of the theory of Bardeen, Cooper and
Schrieffer~\cite{BCS1957}. With $E=(q^2+M_{\rm eff}^2)^{1/2}$ the eigenvalues along the amplitude and along the phase of the
vacuum direction are
\begin{equation}
\lambda^{\rm amp}(\omega)=\frac{G_\SH}{4\pi^2}\int_0^{\Lambda_\LL}\frac{q^2(q^2+m_\LL^2)\,dq}{E\,(E^2-\omega^2/4)},\qquad
\lambda^{\rm ph}(\omega)=\frac{G_\SH}{4\pi^2}\int_0^{\Lambda_\LL}\frac{q^2E\,dq}{E^2-\omega^2/4},
\label{eq:lam}
\end{equation}
as sums over the nodes of the grid with its weights. The Python program finds that its response equals these sums to
{{YH7_DEV}} at five frequencies, in both versions. At $\omega=0$, $\lambda^{\rm ph}(0)=G_\SH I_1(M_{\rm eff})/(4\pi^2)=1$ is the
gap equation and the Goldstone theorem. Since $E^2=q^2+m_\LL^2+(d^0_\SH)^2/4$, the integrands obey, at every momentum,
\begin{equation}
\frac{q^2+m_\LL^2}{E\,(E^2-\omega^2/4)}-\frac1E=\frac{\omega^2-(d^0_\SH)^2}{4E\,(E^2-\omega^2/4)},
\end{equation}
and with the gap equation
\begin{equation}
\lambda^{\rm amp}(\omega)-1=\frac{\omega^2-(d^0_\SH)^2}{4}\,\frac{G_\SH}{4\pi^2}\int_0^{\Lambda_\LL}\frac{q^2\,dq}{E\,(E^2-\omega^2/4)} .
\label{eq:identity}
\end{equation}
Below the threshold every term of the integral is positive, so $\lambda^{\rm amp}=1$ only at
\begin{equation}
\omega^{\rm H}_\SH=d^0_\SH ,
\label{eq:omegaH}
\end{equation}
exactly, for every coupling, mass and cutoff and on every grid with positive weights. The mode is bound by
$2M_{\rm eff}-d^0_\SH=(4m_\LL^2+(d^0_\SH)^2)^{1/2}-d^0_\SH={{BIND2}}$. The Wolfram script derives (\ref{eq:identity}) and the
positivity symbolically and finds the root at $d^0_\SH$ to 22 digits. Result (\ref{eq:omegaH}) is the analogue of the
amplitude mode of a superconductor at twice the gap parameter~\cite{LittlewoodVarma1981,PekkerVarma2015}; here $d^0_\SH/2$
plays the role of the gap parameter, and the mass of the fields raises the lowest energy of a quasi-particle from
$d^0_\SH/2$ to $M_{\rm eff}$, so that the mode is bound.

\emph{The partners of the amplitude mode.} On the eight real components of the four complex pair fields the Python
program finds four eigenvalues equal to 1 at $\omega=0$ and four more at $\omega=d^0_\SH$. The generators $T^a+T^a_R$ of
$SU(2)_V$ commute with the response, and on each set the Casimir $\sum_a(T^a+T^a_R)^2$ has the eigenvalues 0, 2, 2, 2: a
singlet and a triplet. At $\omega=0$ the triplet becomes the longitudinal parts of $W^\pm_\mu$ and $Z_\mu$ and the singlet is
the Goldstone boson of $F^{\rm tot}_\SH$; at $\omega=d^0_\SH$ the singlet is the amplitude mode, and the triplet holds a
neutral scalar and a charged scalar with its antiparticle ($Q=T^3+T^3_R$ on the pair fields). In the first version the
same count gives three Goldstone modes and one amplitude mode, as in the Standard Model.

\subsection{The dependence on the coupling and the mass}

In HS (0,0) the gap equation $G_\SH I_1(M_{\rm eff})/(4\pi^2)=1$ gives $d^0_\SH$ in closed form, so the amplitude mode and the
scale can be followed over the coupling. Table~\ref{tab:scan} gives 32 points: $m_\LL=0.1$ and $0.2$,
$G_\SH/G^{\rm crit}_\SH$ from 1.02 to 3, both versions. In every point the Rust program finds the self-consistent state, the
amplitude mode at $d^0_\SH$ to {{H8D_DEV}}, and $\rho_\SH=1$ in the second version; $\omega^{\rm H}_\SH/v_\SH$ lies between
{{SCAN_MIN2}} and {{SCAN_MAX2}} in the second version, against $m_H/v={{MH_SM_V}}$.

{{T_SCAN}}

%=====================================================================================================================
\section{The dynamics in time}
\label{sec:dynamics}

The time-dependent mean field of method $X$ evolves the quasi-particle vectors of every momentum by the BdG matrix of the
self-energies that the current local matrices produce,
\begin{equation}
i\,\partial_t\,\mathsf O(\vec q,t)=\mathsf H(\vec q;\mathsf N(t),\mathsf P(t))\,\mathsf O(\vec q,t),
\label{eq:tdmf}
\end{equation}
whose linearization about a stationary state is the linear response (\ref{eq:dR}). The program integrates it with the
backward differentiation formulas of CVODE~\cite{SUNDIALS2005,SUNDIALS2022} (relative tolerance $10^{-12}$, absolute
$10^{-14}$, the Newton iteration replaced by a fixed-point iteration), on the same grid as the stationary states. A run
starts from the ground state with the pair self-energy displaced by the relative amount $\epsilon^{\rm kick}$ along the
vacuum direction (an amplitude run) or along the custodial triplet (a triplet run), and records $d(t)$; its spectrum is
the Fourier transform over the run, with the peak frequency $\omega^{\rm pk}$ and the half-width $\Gamma^{\rm pk}$, compared
with the half-width $\Gamma^{\rm win}$ that the length of the run alone gives. Table~\ref{tab:runs} lists the runs with the
conservation of the energy, of the charges and of the orthonormality of the vectors, and Table~\ref{tab:modes} the modes.

{{T_RUNS}}

{{T_MODES}}

In HS (0,1) and HS (0,0) the amplitude run oscillates at $d^0_\SH$ to {{DY_DEV_AMP}}, without damping, and the triplet run
at the same frequency. In HS (1,1) the mode lies at {{DY_RATIO_HS11}}\,$d^0_\SH$, below $2E^{X,{\rm gap}}$, undamped, and the
triplet {{DY_SPLIT_HS11}} per cent above it. In KS and CJT the spectrum peaks at {{DY_AMP_KS}}, above $2E^{X,{\rm gap}}={{DY_GAP2_KS}}$,
{{DY_WR_KS}} times wider than an undamped oscillation: the mode decays into pairs. With a chemical potential $\mu=0.01$ of
the particle number the phase turns at the rate $2\mu$ to {{DY_DEV_JOS}} (the Josephson relation), and a displacement by
$\epsilon^{\rm kick}=0.3$ in HS (0,0) oscillates at {{DY_LARGE_RATIO}}\,$d^0_\SH$, softer than the small oscillation.

%=====================================================================================================================
\section{Toward the gravitational sector: Lord's construction}
\label{sec:gravity}

Lord~\cite{Lord1971} considered the group of the local changes of representation of the two-component spinor algebra,
under which the free field equations of particles of any spin are invariant, and made it a gauge group in the sense of
Yang and Mills. Gauge invariance and Lorentz covariance then require auxiliary fields, which Lord interpreted as the
potentials of a spin-two field; that field is self-interacting, and the Bargmann-Wigner equations of the linear spin-two
field, modified by the self-interaction, give a nonlinear theory identical with Einstein's, the auxiliary fields becoming
in addition the gauge field of the general coordinate transformations. In the present model the factor $SL(2,\mathbb C)$
of the gauge group acts on the spinor indices of the three fields with the connection $\Gamma^\LL_\mu$ of (\ref{eq:cov}),
a pure gauge on flat space-time. The program is to replace $\Gamma^\LL_\mu$ by Lord's gauge field and its auxiliary fields,
to couple the spinor fields and their pairs to it, and to obtain the gravitational field equations with the
energy-momentum tensor of the spinor fields as their source; the route of Utiyama and Kibble~\cite{Utiyama1956,Kibble1961,Hehl1976}
gauges the local Lorentz group in the same spirit. None of this is computed here.

%=====================================================================================================================
\section{The checks}
\label{sec:checks}

The computation was made in stages, each with a program in Rust, one in Python and a script in the Wolfram Language:
stage 2, the group theory and the quantum numbers (sh\_quantum, Sec.~\ref{sec:model}); stage 3, the normal vacuum and the
stationary states (sh\_superfluid, Secs.~\ref{sec:methods} to \ref{sec:states}); stages 4 and 5, the effective potential,
the masses of the gauge bosons and the amplitude mode (sh\_higgs, Secs.~\ref{sec:potential} to \ref{sec:amplitude}); and
stage 6, the dynamics in time (sh\_dynamics, Sec.~\ref{sec:dynamics}). Stage 1 set out the hypotheses.
Every statement of the article and of this Supplemental Material is tested by checks of the programs; each check
compares a computed deviation with a tolerance, or decides a statement that is true or false from the facts that the
program states. The statements of the checks are those that the programs write; where they cite the hypotheses or the
equations of the manuscripts of the stages, the tables below cite instead the place of this Supplemental Material that
states the same thing. The verdicts: the Rust programs {{VERDICTS_RUST}}; the Python programs {{VERDICTS_PYTHON}}; the Wolfram
scripts {{VERDICTS_WOLFRAM}}; the notebook, which repeats the checks of the Wolfram scripts, {{VERDICT_NOTEBOOK}}. The
tables list every check of the Rust and the Python programs with its result, deviation, tolerance and statement.

{{CHECKS_TABLES}}

%=====================================================================================================================
\section{Data and programs}
\label{sec:data}

The repository of PROJECT spacetimes~\cite{Repo}, folder \texttt{Revision\_SM\_Higgs}, holds the programs: the Rust
workspace \texttt{rust/} (the crates \texttt{shquantum}, \texttt{shsuperfluid}, \texttt{shhiggs} and \texttt{shdynamics}, the
library \texttt{sfcore} and the CVODE crate \texttt{sfsundials}), the Python programs \texttt{python/sh\_quantum.py},
\texttt{sh\_superfluid.py}, \texttt{sh\_higgs.py}, \texttt{sh\_dynamics.py} and their figure programs, the Wolfram scripts
\texttt{wolfram/source/sh\_*.wls}, the notebook \texttt{superfluid.nb}, and the build of this article,
\texttt{paper\_prd/scripts/build\_paper.py}. Table~\ref{tab:files} lists the data files and reports that this article reads,
with their SHA-256 checksums.

{{DATA_TABLE}}

\bibliography{references}

\end{document}
